\section{Debate Generation Pipeline Details}
\label{app:pipeline}

This appendix expands Section~\ref{sec:pipeline}. The pipeline is shared
identically between baseline and experimental runs; all quality differences
reported in the main body therefore come from the speech generation model
and training procedure, not from asymmetric preparation.

\subsection{Belief Tree Structure}

From a resolution $R$, the pipeline constructs a typed belief tree in three
passes.

\paragraph{Values.}
A DSPy \textsc{ValueExtractor} module generates 4 topic-specific values,
2 per side, phrased as terminal goods the side should be understood to
defend (e.g., \emph{economic mobility}, \emph{institutional trust}).
Values are the roots of the belief tree.

\paragraph{Beliefs.}
A \textsc{BeliefGenerator} expands each value into 2 child beliefs per
side, then refines each to depth-2, yielding approximately 32 beliefs
total and 16 arguments per side. Each belief $b$ carries a Bayesian
credence $c(b)\in[0,1]$; the prior is set by the generator and is later
updated as posteriors conditioned on speeches delivered.

\paragraph{Arguments.}
Leaf beliefs are instantiated as typed arguments drawn from standard
debate theory: \textsc{uniqueness}, \textsc{link}, \textsc{impact},
\textsc{solvency}, and \textsc{turn}. Typing is important: it gives the
downstream speech stages structured slots (e.g., an impact argument must
identify magnitude, timeframe, probability), and it allows the hierarchical
reward (Section~\ref{sec:reward_hierarchy}) to reason about the debate at
the level of argument roles rather than raw sentences.

\subsection{Research Pipeline: v1 vs.\ v2}

\paragraph{v1 (9 LLM calls per hop).}
The original research agent made separate LLM calls for: query
formulation, query rewriting, relevance filtering, per-document
summarization, per-sentence judgement, selection ranking, card tagging,
citation extraction, and hop termination. Typical claims took 3--4 hops,
burning $\sim$30 calls per claim and dominating preparation cost.

\paragraph{v2 (2 LLM calls per hop, consolidated).}
The consolidated loop per claim is:
\begin{enumerate}
    \item \textbf{Query generation (LLM call 1):} A single structured call
    emits $k$ retrieval queries plus termination heuristics.
    \item \textbf{Parallel retrieval:} Weaviate (dense) and Tavily (web)
    run in parallel on each query.
    \item \textbf{Deterministic processing:} Text is cleaned, chunked into
    sentences, and scored with hybrid
    $s = 0.7\cdot \cos(\mathbf{q},\mathbf{s}) + 0.3\cdot\mathrm{BM25}(q,s)$.
    Top sentences are expanded into windows to preserve local context.
    \item \textbf{Selection (LLM call 2):} A single structured call
    consumes the candidate windows and emits the final evidence card via
    Pydantic validation, including bolded selections and sentence IDs.
\end{enumerate}
Most claims terminate after one hop. Total calls per claim drop from
$\sim$30 to $\sim$2--4, a 10$\times$ reduction that is what makes the
downstream branching tree (Section~\ref{sec:branching}) economically
feasible.

\subsection{Evidence Card Schema}

Each card is a Pydantic record with the following fields:
\begin{itemize}
    \item \texttt{tag} --- one-line argumentative claim.
    \item \texttt{cite}, \texttt{fullcite} --- short and full citations.
    \item \texttt{full\_text} --- the complete retrieved passage.
    \item \texttt{card} --- the passage with \textbf{bolded} sentences
    indicating selected material.
    \item \texttt{selected\_texts} --- list of selected sentences as
    strings.
    \item \texttt{selected\_sentence\_ids} --- indices into
    \texttt{full\_text}, the \emph{only} handle the speech generator can
    cite.
    \item \texttt{window\_ranges} --- windows around each selection used
    for context.
    \item \texttt{source\_url} --- retrieval source for audit.
    \item \texttt{relevance\_score} --- retrieval score used by sorting.
\end{itemize}
Because downstream speech generation cites by \texttt{selected\_sentence\_id}
only, the quotation text is mechanically assembled from \texttt{full\_text}
and cannot be fabricated by the model. This is the structural property that
makes evidence verification in Section~\ref{sec:reward_hierarchy} a
sentence-ID check rather than a natural-language fidelity judgment.

\subsection{Perspective Building}

Given the completed belief tree, a \textsc{PerspectiveBuilder} compiles
side-specific prompt contexts:
\begin{itemize}
    \item \textbf{Constructives} receive a pruned depth-1 view: 2 of the
    4 root beliefs for the side, chosen to cover complementary ground.
    This forces the constructive speech to commit to a narrow thesis
    rather than diffuse across all available arguments.
    \item \textbf{Rebuttals} receive the full sub-belief tree in SEAM
    form (Setup, Evidence, Analysis, Move), which aligns naturally with
    rebuttal block structure in competitive debate.
    \item \textbf{CX} receives the opponent's most recent
    constructive along with attackable beliefs.
\end{itemize}
Different branches of the debate tree (Section~\ref{sec:branching}) are
seeded with different D1 belief combinations, producing natural debate
diversity without any hand-written prompts or per-topic tuning.

\subsection{Why Automation Matters}

\paragraph{Reproducibility.}
Every piece of training data is reproducible from the original resolution
plus the research index snapshot: there is no human-curated brief, no
pre-written argument, no case file. This matters for scientific validity;
it also matters for operational robustness when we replay or audit
training runs.

\paragraph{Scale.}
The pipeline runs end-to-end on a resolution in minutes, making tournament
generation (68 debates) and training-scale generation (thousands of
debates) feasible on a small budget ($\sim$\$750--1{,}000 total;
Appendix~\ref{app:costs}).

\paragraph{Baseline/experimental comparability.}
The same pipeline produces the inputs for both the base Qwen3-30B-A3B and
all trained checkpoints. A difference in downstream win rate therefore
cannot be attributed to a richer prep stage, a better research index, or
stronger evidence --- the substrate is identical, so the delta isolates the
generation model. This is a property we rely on for the tournament
analysis in Section~\ref{sec:analysis}.

\paragraph{Ablation surface.}
Because the pipeline is structured rather than monolithic, components can
be ablated independently: belief tree depth, research-hop limits, SEAM
structure, and the four-stage decomposition itself. Each is a natural
axis of ablation and a natural training target, independent of the
specific model being trained.
